\documentclass[a4paper, 11pt]{article}

\usepackage[T1,T2A]{fontenc}
\usepackage[utf8]{inputenc}
\usepackage[english, russian]{babel}

\usepackage{comment}
\usepackage{fullpage}
\usepackage{amssymb}
\usepackage{amsmath}
\usepackage{amsthm}
\usepackage{enumitem}

% %\DeclareMathOperator{\ord}{ord}
\DeclareMathOperator{\chr}{char}
\DeclareMathOperator{\Mod}{mod}
\DeclareMathOperator{\Gal}{Gal}
\DeclareMathOperator{\N}{N}
\DeclareMathOperator{\Tr}{Tr}
\DeclareMathOperator{\SL}{SL}

\begin{document}

\begin{center} \Large \bf Листок 8\\ Тема 8 (3.1). Характеры. Суммы Гаусса \end{center}

\noindent\textbf{Упражнения и задачи}

\begin{enumerate}[topsep=0pt]
    \item Пусть $G$ --- конечная абелева группа, $H$ --- собственная подгруппа, $g\in G$, $g \not\in H$. Докажите, что существует характер $\chi$ группы $G$ такой, что $\chi(g)\neq 1$ и $\forall h\in H$ $\chi(h)=1$. %[LN Ex 5.1]
    \item Пусть $G$ --- конечная абелева группа, $m>0$ --- целое. Докажите, что $g\in G$ является $m$-ой степенью в $G$ $\iff$ для всех характеров $\chi$ группы $G$ с $\chi^m=\chi_0$ выполняется $\chi(g)=1$. %???
    \item Докажите, что для аддитивного характера $\psi=\psi_1$, $\psi_1(\alpha)=e^{2\pi i\,\mathrm{Tr}(\alpha)/p}$, поля $\mathbb{F}_q/\mathbb{F}_p$ (где $\mathrm{Tr}=\mathrm{Tr}_{\mathbb{F}_q/\mathbb{F}_p}$; при $q=p$ это $\psi(x)=e^{2\pi i x/p}$) для всех $\alpha\in\mathbb{F}_q$, $j\geqslant 0$ справедливо $\psi_1(\alpha^{p^j})=\psi_1(\alpha)$. %[LN Ex 5.7]
    \item Пусть $\chi'$ --- мультипликативный характер $\mathbb{F}_{q^s}$ порядка $m$, $\chi$ --- ограничение $\chi'$ на $\mathbb{F}_q$. Докажите, что $\chi$ --- мультипликативный характер $\mathbb{F}_q$ порядка $m/(m,(q^s-1)/(q-1))$. %[LN Ex 5.8]
    \item Пусть $\chi$ --- мультипликативный характер $\mathbb{F}_q$, $\chi'=\chi\circ N_{\mathbb{F}_{q^s}/\mathbb{F}_q}$ --- его поднятие на $\mathbb{F}_{q^s}$. Докажите, что $\chi'$ --- мультипликативный характер $\mathbb{F}_{q^s}$ и $\chi'(a)=\chi(a)^s$ $\forall a \in\mathbb{F}_q^{*}$. %[LN Ex 5.10]
    \item Пусть $H$ --- множество функций $f:\mathbb{F}_p\to\mathbb{C}$, определим скалярное произведение $(f,g)=\frac{1}{p}\sum_{x\in\mathbb{F}_p}f(x)\overline{g(x)}$. Пусть $\psi_a(x)=e^{2\pi i ax/p}$, $a,x\in\mathbb{F}_p$. Докажите, что $(\psi_a)_{a\in\mathbb{F}_p}$ --- ортонормированный базис $H$.
    \item Пусть $f:\mathbb{F}_q \rightarrow \mathbb{C}$, $\psi\neq\psi_0$ --- аддитивный характер $\mathbb{F}_q$, $\hat{f}(s) = \frac{1}{q}\sum_{u\in\mathbb{F}_q} f(u)\overline{\psi(su)}$ --- конечное преобразование Фурье. Докажите, что $f(t)=\sum_{s\in\mathbb{F}_q} \hat{f}(s) \psi(st)$. %[IR Ex 10.23]
    \item Пусть $p\neq 2$, $ab\not\equiv 0\,(p)$, $\chi$ --- квадратичный характер $\mathbb{F}_p^*$, $\psi_a(x)=e^{2\pi i ax/p}$, $G(\chi,\psi_a)=\sum_{x\in\mathbb{F}_p^*}\chi(x)\psi_a(x)$. Докажите, что
    \begin{itemize}[topsep=0pt]
        \item $G(\chi,\psi_a)G(\chi,\psi_b) = \left(\dfrac{-ab}{p}\right) p$; %[БШ стр 26 упр 5]
        \item $\sum_{a\in\mathbb{F}_p} G(\chi,\psi_a)=0$. %[БШ стр 26 упр 6]
    \end{itemize}
    \item Докажите свойства обобщенных сумм Гаусса для конечного поля $\mathbb{F}_q/\mathbb{F}_p$. Пусть $\chi$ --- мультипликативный, $\psi$ --- аддитивный характер $\mathbb{F}_q$, $\psi_b(x)=\psi(bx)$, $\bar\psi(x)=\psi(-x)$, $G(\chi,\psi)=\sum_{x\in\mathbb{F}_q^*}\chi(x)\psi(x)$.
    \begin{itemize}[topsep=0pt]
        \item $G(\chi,\psi_{ab})=\bar{\chi}(a) G(\chi, \psi_b)$, $a\in\mathbb{F}_q^*$, $b\in \mathbb{F}_q$;
        \item $G(\chi,\bar{\psi}) = \chi(-1) G(\chi,\psi)$;
        \item $G(\bar{\chi},\psi) = \chi(-1) \overline{G(\chi,\psi)}$;
        \item $G(\chi,\psi)G(\bar{\chi},\psi)=\chi(-1) q$, $\chi\neq\chi_0$, $\psi\neq\psi_0$;
        \item $G(\chi^p,\psi_b)=G(\chi,\psi_{\sigma(b)})$ при $\psi=\psi_1$, $b\in\mathbb{F}_q$, $\sigma$ --- автоморфизм Фробениуса;
        \item $G(\chi_0,\psi_0)=q-1$, $G(\chi_0,\psi)=-1$ при $\psi\neq\psi_0$, $G(\chi,\psi_0)=0$ при $\chi\neq\chi_0$.
    \end{itemize} %[LN Th 5.12]
\end{enumerate}

% \noindent\textbf{SageMath}
% \begin{itemize}[topsep=0pt]
%     \item Исследуйте основные функции SageMath связанные с группой характеров конечных абелевых групп:
%     \begin{itemize}[noitemsep,topsep=0pt]
%         \item \texttt{character\_table()}.
%      \end{itemize}
% \end{itemize}

\noindent\textbf{Темы для самостоятельного изучения}
\begin{itemize}[topsep=0pt]
    \item Доказательство соотношения Хассе--Дэвенпорта. [LN], \S 5.2.
    \item Доказательство квадратичного закона взаимности через суммы Гаусса. [IR], \S 7.3; [LN], \S 5.2.
\end{itemize}


\end{document}
